\section{Discussion}
\label{sec:discussion}

% §8.1 What this paper does not claim
%   - Verification does not solve alignment. Specification and
%     generalization remain open problems.
%   - Mesa-optimization does not reduce to verification. An agent that
%     has internalized the wrong objective will pursue it faithfully
%     regardless of how well verification works.
%   - Crypto is not the only way to exogenize verification. Hardware
%     enclaves, physical separation, and human-in-the-loop are all
%     forms of substrate exclusivity. Crypto is the form that scales
%     without requiring continuous human attention.

% §8.2 Limitations
%   - The commitment protocol assumes the capability poset is
%     well-defined enough that "capability c_k" has a testable meaning.
%     For vaguely defined capabilities, the exercise protocol may not
%     be constructible.
%   - Substrate exclusivity is a physical assumption. If a digital
%     agent gains write access to a biological substrate (e.g., through
%     brain-computer interfaces), the substrate barrier degrades.
%   - The ledger's append-only property requires at least one honest
%     node per substrate. If all nodes on a substrate are compromised,
%     the substrate's contribution to consensus is lost.
%   - Zero-knowledge capability proofs require the capability to be
%     formalizable as a predicate on the poset. Some capabilities may
%     resist formalization.

% §8.3 Open problems
%   - Exercise protocol design: what constitutes a valid test of
%     capability c_k? This is domain-specific and may require
%     capability-specific protocols.
%   - Partial verification: can we verify that an agent has "at least"
%     capability c_k without testing the exact boundary?
%   - Dynamic capabilities: capabilities change over time. The ledger
%     records point-in-time verifications; continuous verification
%     requires periodic re-commitment.
%   - Cross-substrate witness incentives: why should a witness on
%     substrate s' spend resources verifying claims by agents on
%     substrate s? The vol_P-maximization incentive provides a
%     structural answer (verification is a public good that benefits
%     the whole population), but the free-rider problem is not
%     formally resolved.

% TODO: write full discussion.
